%% THERMODYNAMIQUE %%
%
\section{Direction et vitesse du vent � partir des composantes}

\itemstyle{Op\'erateur~:} \operfmt{UV2ddff}
\index{UV2ddff}

%%%%%OBJET%%%%%
\itemstyle{Objet~: }
Calcul des composantes polaires du vent (direction, vitesse), selon la
convention m�t�orologique usuelle, � partir de composantes cart�siennes
dans le rep�re g�ographique
local. La convention m�t�orologique consiste � indiquer la direction
d'o� vient le vent, exprim�e dans le sens r�trograde � partir du nord
g�ographique.

%%%%%ARGUMENTS%%%%%
\bigskip
%%%%%ENTREES%%%%
\begin{tabarg}
\multicolumn{3}{l}{\titretab{Entr\'ees}}\\
\hline
$U$ & Vecteur, \texttt{float} & Composante Ouest-Est du vent ($m/s$)\\
$V$ & Vecteur, \texttt{float} & Composante Sud-Nord du vent ($m/s$) \\	
\end{tabarg}

\medskip
%%%%%SORTIES%%%%
\begin{tabarg}
\multicolumn{3}{l}{\titretab{Sorties}}\\
\hline
$DD$ & Vecteur, \texttt{float} & Direction g\'eographique du vent selon convention m\'et\'eorologique, dans l'intervalle $[0, 360[$ (\textit{degr\'es})\\
$FF$ & Vecteur, \texttt{float} & Vitesse du vent (m\^eme unit\'e que $U$ et $V$)\\
\end{tabarg}


%%%%%FORMULE%%%%
\itemstyle{Formule}

%% Formule de l'operateur : UV2ddff
%% Dans ce fichier, les �ventuels accents sont cod�s
%% en ISO-8859-15 (aka latin9)
%% La ligne suivante donne trois exemples 
%% e accent aigu : �, e accent grave : �, e-dans-l'o : �

$FF = \sqrt{ U^2 + V^2 } $ \newline

Si $V=0$ alors  \newline
\begin{tabular}{p{3cm}|l}
 & si $U=0$ alors $DD=0$ \\
 & si $U > 0$ alors $DD=-\pi /2$ \\
 & si $U < 0$ alors $DD=\pi /2$ \\
\end{tabular}

Si $V > 0$ alors \newline
\begin{tabular}{p{3cm}|l}
 & $DD = \arctan ( U / V)  + \pi$ \\
\end{tabular}

Si $V < 0$ alors \newline
\begin{tabular}{p{3cm}|l}
 & $DD = \arctan ( U / V)$  \\
\end{tabular}
\newline
Les �ventuelles directions n�gatives sont rendues positives : si $DD < 0$ alors
$ DD = DD + 2*\pi $

Puis conversion en degr�s :
$ DD = DD * \frac{180}{\pi}$



%%%%%COMMENTAIRE%%%%
\itemstyle{Commentaire~: }

Bien que le r�sultat pour la direction s'appelle $DD$, comme dans les
descriptions du codage des messages m�t�o, il est exprim� comme un nombre r�el,
et pas selon ces r�gles de codage (dans lesquelles la valeur est arrondie �
la dizaine de degr�s la plus proche, dans l'intervalle $[10, 360]$, $0$ �tant
r�serv� au vent nul).

Bien que le r�sultat pour la vitesse s'appelle $FF$, cette grandeur est bien
une vitesse (distance parcourue pendant un certain temps), et pas une force
(qui serait exprim�e en Newton).

%%%%%AUTEUR%%%%%
\itemstyle{Auteur~:}
CNRM/GMEI/TRAMM

%%%%%REFERENCES%%%%%
\itemstyle{R\'ef\'erences~:}
